\paragraph{可见 strictification 塔的有限饱和、异常坐标计数与终端普适商}
在有限阿贝尔 lift 群口径下，定理 \ref{thm:pom-bc-visible-quotient-eventual-stability} 已表明可见 strictification 核链 \(\{N_m^{\mathrm{BC}}\}_{m\ge 1}\) 在有限阶段后稳定。把这一有限步饱和与 Pontryagin 对偶分类结合起来，便可把“潜在异常方向”的数量与最终可见商的普适性完全压缩为两个显式的有限群不变量。

\begin{theorem}[独立异常坐标的精确计数律]\label{thm:xi-terminal-gbc-anomaly-coordinate-count}
\leanverified{paper\_xi\_terminal\_gbc\_anomaly\_coordinate\_count}
沿用定义 \ref{def:pom-bc-quotient} 的记号，并假设 \(G\) 为有限阿贝尔群。对每个分辨率 \(m\)，定义独立异常坐标群
\[
\mathcal A_m:=\widehat G/\widehat{G_m^{\mathrm{BC}}}.
\]
则存在自然同构
\[
\mathcal A_m\cong \widehat{N_m^{\mathrm{BC}}},
\]
因而其非平凡独立坐标的数量恰为
\[
\#\mathcal A_m^{\mathrm{nz}}
:=
|\mathcal A_m|-1
=
|N_m^{\mathrm{BC}}|-1.
\]
此外，\(\#\mathcal A_m^{\mathrm{nz}}\) 随 \(m\) 单调不增；并且若
\[
N_\infty^{\mathrm{BC}}
:=
\bigcap_{m\ge 1}N_m^{\mathrm{BC}},
\]
则在定理 \ref{thm:pom-bc-visible-quotient-eventual-stability} 的饱和阶段之后，有
\[
\#\mathcal A_m^{\mathrm{nz}}=|N_\infty^{\mathrm{BC}}|-1.
\]
因此，潜在异常方向在有限分辨率上的全部独立坐标，并不是抽象余维，而是由 strictification 核的大小精确计数。
\end{theorem}

\begin{proof}
由推论 \ref{cor:pom-bc-pontryagin-dual-classification}，
\[
\widehat G/\widehat{G_m^{\mathrm{BC}}}
\cong
\widehat{N_m^{\mathrm{BC}}},
\]
即得所示自然同构。
又因有限阿贝尔群与其 Pontryagin 对偶等势，故
\[
|\mathcal A_m|
=
|\widehat{N_m^{\mathrm{BC}}}|
=
|N_m^{\mathrm{BC}}|.
\]
去掉平凡角色即得
\[
\#\mathcal A_m^{\mathrm{nz}}=|N_m^{\mathrm{BC}}|-1.
\]

再由定理 \ref{thm:pom-bc-visible-quotient-eventual-stability}，
\[
N_{m+1}^{\mathrm{BC}}\le N_m^{\mathrm{BC}},
\]
故 \(|N_m^{\mathrm{BC}}|\) 随 \(m\) 单调不增，从而 \(\#\mathcal A_m^{\mathrm{nz}}\) 亦单调不增。并且在某个有限阈值 \(m_0\) 之后，
\[
N_m^{\mathrm{BC}}=N_\infty^{\mathrm{BC}}\qquad(m\ge m_0),
\]
于是
\[
\#\mathcal A_m^{\mathrm{nz}}=|N_\infty^{\mathrm{BC}}|-1
\qquad(m\ge m_0).
\]
证毕。
\end{proof}

\begin{theorem}[稳定 strictification 商的终对象性质]\label{thm:xi-terminal-gbc-stabilized-terminal-object}
\leanverified{paper\_xi\_terminal\_gbc\_stabilized\_terminal\_object}
仍设 \(G\) 为有限阿贝尔群。考虑范畴 \(\mathsf{Str}^{\mathrm{vis}}_\infty(G)\)，其对象为满射
\[
\pi:G\twoheadrightarrow H
\]
满足存在某个阈值 \(m_\pi\)，使得对一切 \(m\ge m_\pi\)，由 \(\pi\) 诱导的可见扭曲通道在分辨率 \(m\) 上均已严格交换；态射为使商映射交换的群同态。则稳定 strictification 商
\[
G_\infty^{\mathrm{BC}}
\cong
G/N_\infty^{\mathrm{BC}}
\]
是范畴 \(\mathsf{Str}^{\mathrm{vis}}_\infty(G)\) 中的终对象。

换言之，任何在足够高分辨率上最终 strict 的可见群商，都且仅都经由 \(G_\infty^{\mathrm{BC}}\) 唯一因子化。
\end{theorem}

\begin{proof}
任取对象 \(\pi:G\twoheadrightarrow H\in\mathsf{Str}^{\mathrm{vis}}_\infty(G)\)。按定义，存在 \(m_\pi\) 使得对所有 \(m\ge m_\pi\)，\(\pi\) 在分辨率 \(m\) 上均给出严格交换的可见商。

另一方面，由定理 \ref{thm:pom-bc-visible-quotient-eventual-stability}，存在 \(m_0\) 使得
\[
G_m^{\mathrm{BC}}\cong G_\infty^{\mathrm{BC}}
\qquad(m\ge m_0).
\]
令 \(M:=\max\{m_\pi,m_0\}\)。则 \(\pi\) 在分辨率 \(M\) 上已严格交换。由命题 \ref{prop:pom-bc-quotient-universal} 的普适性，\(\pi\) 必唯一因子化为
\[
G\twoheadrightarrow G_M^{\mathrm{BC}}\longrightarrow H.
\]
又因 \(M\ge m_0\)，有
\[
G_M^{\mathrm{BC}}\cong G_\infty^{\mathrm{BC}},
\]
故 \(\pi\) 亦唯一因子化为
\[
G\twoheadrightarrow G_\infty^{\mathrm{BC}}\longrightarrow H.
\]
这正是终对象的定义。证毕。
\end{proof}

\endinput
